% Generated by tools/edn_latex.py from reports/edn.md.
% Do not edit: change reports/edn.md and run `make edn`.
\ednentry{
  id = {EDN-73},
  title = {\texorpdfstring{\texttt{learning\_\allowbreak{}rate}}{learning\_rate} and \texorpdfstring{\texttt{num\_\allowbreak{}leaves}}{num\_leaves} are not tuned, and any later tuning follows a written protocol},
  note = {\textbf{Extended by EDN-78} (2026-10-06, issue \#88): the protocol below now covers both LightGBM variants and \texttt{min\_\allowbreak{}child\_\allowbreak{}samples}, its decision step is \texttt{recommenditos/\allowbreak{}modeling/\allowbreak{}tune.\allowbreak{}py}, and every fit of a search is measured and tracked.
This entry is kept as the record of the first sweep and of the protocol as it was written.},
  date = {2026-10-05},
  milestone = {M3: Quality Assurance},
  topic = {Hyperparameter Tuning, Evaluation Protocol, Experiment Tracking},
  participants = {@lukas2510, decided by the repository owner},
  decision = {\texttt{learning\_\allowbreak{}rate: 0.\allowbreak{}05} and \texttt{num\_\allowbreak{}leaves: 63} stay as they are for both LightGBM variants, because a bounded sweep around them found no point the validation split can tell apart from them.
The sweep is kept as the evidence that this region is flat, and any later tuning of a committed hyperparameter follows a protocol written down in \texttt{docs/\allowbreak{}docs/\allowbreak{}pipeline.\allowbreak{}md}, ``Tuning a hyperparameter'':
queued \texttt{dvc exp run -\allowbreak{}-\allowbreak{}queue -\allowbreak{}S .\allowbreak{}.\allowbreak{}.\allowbreak{} train@\allowbreak{}<variant>} experiments screen a sweep, so \texttt{evaluate} never runs and the test split takes no part in the choice, and they are tracked in their own MLflow experiment, \texttt{recommenditos-\allowbreak{}price-\allowbreak{}tuning}, so the ladder's \texttt{recommenditos-\allowbreak{}price} holds only pipeline runs;
\texttt{reports/\allowbreak{}analysis/\allowbreak{}tuning\_\allowbreak{}sweep.\allowbreak{}py} then re-fits the shortlist on the validation split and decides by \texttt{validation\_\allowbreak{}l1\_\allowbreak{}log\_\allowbreak{}price}, replacing a committed value only when its simultaneous 95\,\% interval lies below zero -- a paired bootstrap of the validation L1 difference, resampled by seller group, widened by the max-statistic bootstrap so that the intervals of all the sweep's points hold together.},
  rationale = {\emph{Alternatives considered:}
\begin{itemize}[leftmargin=*, topsep=2pt]
\item \textbf{Option A (chosen): do not tune now, and write the protocol down.}
\newline Pros: no number moves for a gain the data cannot resolve, so the report's numbers are final for this delivery.
It avoids the winner's curse: of eleven noisy comparisons against the committed point, the best-looking one is the one most likely to look better than it is, and a simultaneous interval is what keeps the protocol from adopting it.
It costs no compute, and the protocol makes the next tuning attempt cheap and honest instead of ad hoc.
\newline Cons: the model is untuned in the conventional sense, and the evidence covers a small grid of two knobs for one variant; \texttt{min\_\allowbreak{}child\_\allowbreak{}samples}, \texttt{feature\_\allowbreak{}fraction} and L2 regularisation are not parameters yet, and \texttt{lgbm-\allowbreak{}extended} was not swept.
\item \textbf{Option B: adopt the validation-L1 winner, \texttt{learning\_\allowbreak{}rate} 0.025 with \texttt{num\_\allowbreak{}leaves} 63, under the same protocol.}
\newline Pros: the lowest validation L1 of the grid, and a model that can be called tuned.
\newline Cons: it would fail the protocol it is adopted under.
The gain is inside the noise even point by point: L1 -0.00031 with a per-point 95\,\% interval of [-0.00086, +0.00025], better in 86\,\% of the draws, and MdAPE -0.009 pp with an interval of [-0.11, +0.09] pp; its simultaneous L1 interval is [-0.00109, +0.00047].
It costs 1.55x the trees and so the fit time and energy, a booster of about 12 MB and about 41 ms for one row's SHAP export, and it moves every number in the report for nothing measurable.
\item \textbf{Option C: a wider search, random or Optuna over more knobs, selected by seller-grouped k-fold cross-validation on train plus validation.}
\newline Pros: cross-validation is a less noisy selection than one validation split, and a larger space might hold a real gain.
\newline Cons: new code and new parameters, about five times the fits per point, more energy, and a departure from the problem specification's ``the validation set is used for early stopping and tuning'' -- all for a gain the flat grid suggests is small.
Left open for a later ticket if a criterion ever comes to depend on a fraction of a point.
\item \textbf{Option D: sweep with \texttt{dvc exp} over the whole pipeline and pick by the test metrics \texttt{dvc exp show} prints.}
\newline Pros: the least setup, and the metric is the one the gate reads.
\newline Cons: rejected, because it leaks the test split into the model the gate then judges, which makes SC-01 to SC-04 an optimistic estimate of a choice the test rows were part of.
\end{itemize}
\par
\emph{Why:} \texttt{reports/\allowbreak{}analysis/\allowbreak{}tuning\_\allowbreak{}sweep.\allowbreak{}py} fitted \texttt{lgbm-\allowbreak{}basic}, the candidate (EDN-62), at \texttt{learning\_\allowbreak{}rate} 0.025, 0.05 and 0.1 against \texttt{num\_\allowbreak{}leaves} 31, 63, 127 and 255, with the ceiling at 20,000 so early stopping decided every point (it did, at best rounds from 217 to 4,066), through the stage's own \texttt{fit\_\allowbreak{}variant} on the train and validation matrices only.
Each point's difference from the committed one was bootstrapped 2,000 times, paired and resampled by the 2,784 seller groups of the 8,859 validation rows, with the same draws for every point.
The region is flat: validation MdAPE spans 6.88\,\% to 7.04\,\% across all twelve points, every MdAPE interval spans zero and is about \ensuremath{\pm}0.13 pp wide, and no point improves validation L1 with an interval below zero.
Point by point, five are measurably \emph{worse} in L1: all four at \texttt{learning\_\allowbreak{}rate} 0.1, and 127 leaves at 0.05.
Simultaneously over the eleven, only 127 and 255 leaves at \texttt{learning\_\allowbreak{}rate} 0.1 remain worse, and no point is better.
The two metrics do not even agree on the best point -- \texttt{(0.\allowbreak{}025, 63)} by L1, \texttt{(0.\allowbreak{}05, 255)} by MdAPE -- which is the signature of a choice among noise.
\par
The protocol's parts each follow from the measurement.
Validation L1 is the selection metric because it is what early stopping reads and because its intervals are about 2.5 times narrower than MdAPE's relative to the value, so it resolves differences MdAPE cannot; MdAPE stays the gate's metric, read once on test after the choice.
The interval is simultaneous because a per-point one is not safe across a sweep: with eleven points no better than the committed one, the luckiest clears a per-point 95\,\% interval in about 24\,\% of sweeps if the comparisons were independent, and in at most 27.5\,\% by the union bound.
The max-statistic bootstrap reuses the paired, seller-grouped draws: each point's difference is standardised by its own bootstrap spread, the largest standardised deviation across the points is taken per draw, and its 95th percentile, 2.80 on this grid against 1.96 for one interval, sets every point's half-width.
\texttt{dvc exp run -\allowbreak{}-\allowbreak{}queue} targeted at \texttt{train@\allowbreak{}<variant>} was run once on 2026-10-05 to check the protocol is executable as written: only \texttt{train@\allowbreak{}lgbm-\allowbreak{}basic} ran, the stages before it were skipped as unchanged, \texttt{evaluate} was never reached, and the queue worker inherited the exported tracking variables, which matters because the queued workspace is built from the commit and does not see the gitignored \texttt{.\allowbreak{}env}.
It also showed one trap the protocol warns about: since \texttt{evaluate} does not run, every experiment carries its parent's \texttt{metrics.\allowbreak{}json}, so the test metrics \texttt{dvc exp show} prints belong to the committed model, not to the point.
So the queued runs are for screening, read in MLflow, and the decision is the analysis script's: it re-fits the shortlist on the validation split rather than reading the screening runs' bundles back, which MLflow (under \texttt{model/\allowbreak{}}) and the experiments' DVC outputs do keep, because a re-fit far past the ceiling gives every point its early-stopped model -- the same trees where the screening run early-stopped, and the uncut model where the committed ceiling bound it.
The script covers \texttt{learning\_\allowbreak{}rate} and \texttt{num\_\allowbreak{}leaves} only, and always fits the committed point, whether or not the grid names it.
\href{https://github.com/mlops-2627q1-mds-upc/MLOPS_Recommenditos/issues/88}{\#88} moved its statistics into \texttt{recommenditos/\allowbreak{}modeling/\allowbreak{}tune.\allowbreak{}py}, which runs the decision step for any grid and records every fit, and extended the protocol to \texttt{lgbm-\allowbreak{}extended} and \texttt{min\_\allowbreak{}child\_\allowbreak{}samples} (EDN-78).},
  ai = {Information seeking, Alternative generation, Alternative assessment, Recommendation, Solution generation},
  response = {Accepted with modifications},
  assessment = {AI designed and ran the bounded sweep, added the paired, seller-grouped bootstrap so the question became ``can the validation split tell these apart'' rather than ``which number is lowest'', and extended the grid to \texttt{learning\_\allowbreak{}rate} 0.025 after the first run showed the lower learning rate was the only direction not clearly worse.
It laid out options A to D, recommended A with the protocol, and stopped for the decision rather than committing a tuned value; Lukas chose A as proposed.
AI then ran one queued \texttt{dvc exp} experiment to confirm the documented commands do what the protocol claims before writing them into the pipeline documentation, which is how the \texttt{metrics.\allowbreak{}json} trap and the \texttt{.\allowbreak{}env} point were found, and wrote the decision step into the analysis script.
An independent review of the pull request, also by AI, corrected the first version of the protocol in two places.
Its rule used a per-point interval and called that safe against the winner's curse, which it is not across eleven comparisons.
And it left implicit that the decision step re-fits the shortlist rather than reading the queued runs back, while wrongly saying that neither DVC nor MLflow keeps what that step needs; MLflow keeps each run's bundle under \texttt{model/\allowbreak{}}, and each experiment's DVC outputs keep \texttt{models/\allowbreak{}}.
The rule became simultaneous, the sweep was re-run with the max-statistic intervals, and the conclusion did not change: a simultaneous interval is only wider, so ``no point distinguishable'' holds a fortiori.},
  aievidence = {Claude Code session on 2026-10-05 implementing issue \#64: the sweep's output is \texttt{reports/\allowbreak{}analysis/\allowbreak{}tuning\_\allowbreak{}sweep\_\allowbreak{}results.\allowbreak{}txt}; the options and the recommendation were reported to Lukas with those numbers, and he chose A.},
  evidence = {\href{https://github.com/mlops-2627q1-mds-upc/MLOPS_Recommenditos/issues/64}{issue \#64}; \href{https://github.com/mlops-2627q1-mds-upc/MLOPS_Recommenditos/pull/77}{PR \#77}; \href{https://github.com/mlops-2627q1-mds-upc/MLOPS_Recommenditos/blob/main/reports/analysis/tuning_sweep.py}{\texttt{reports/\allowbreak{}analysis/\allowbreak{}tuning\_\allowbreak{}sweep.\allowbreak{}py}} and its results file (2026-10-05); \href{https://github.com/mlops-2627q1-mds-upc/MLOPS_Recommenditos/blob/main/docs/docs/pipeline.md}{pipeline docs}, ``Tuning a hyperparameter''; \href{https://github.com/mlops-2627q1-mds-upc/MLOPS_Recommenditos/blob/main/docs/docs/model-card.md}{model card}, Hyperparameters as trained; \href{https://github.com/mlops-2627q1-mds-upc/MLOPS_Recommenditos/blob/main/params.yaml}{\texttt{params.\allowbreak{}yaml}}, \texttt{train.\allowbreak{}variants.\allowbreak{}lgbm-\allowbreak{}basic}; \href{https://github.com/mlops-2627q1-mds-upc/MLOPS_Recommenditos/blob/main/docs/docs/problem-spec.md}{problem specification} section 5; \hyperref[EDN-70]{EDN-70}, \hyperref[EDN-14]{EDN-14}.},
}
